% ==========================================================================
% GENERATED by scripts/split_paper.py. Do not edit this file.
%
% Source: imported_papers/Koopman-JEPA-TMLR-rewrite/main.tex
% Part 4: A Controlled Comparison, and How It Is Read
%
% Edits belong upstream, in the published paper. To change what the
% chapter includes, change imports/bounded.toml and re-run
% the splitter.
% ==========================================================================
\section{A Controlled Comparison, and How It Is Read}
\label{sec:arch}

\begin{figure*}[t]
    \centering
    \includegraphics[width=0.76\textwidth]{paper_figures/koopman.png}
    \caption{The four dynamics modules differ only in how the spectrum of the
    transition operator is constrained. \textit{(Left)} An online encoder $\phi$
    maps the observation $\x_t$ to a latent $\z_t$, a dynamics module predicts
    future latents from an action sequence, and the prediction loss compares
    them against targets from an EMA-updated target encoder $\bar\phi$. No
    pixels are decoded during training. \textit{(Right)} The four modules. The
    MLP baseline has no transition matrix and so no spectrum, and the panel
    marks unbounded local amplification instead. Koopman~(Unconstrained) places
    no constraint on the eigenvalues of $\K$, Koopman~(Soft) penalizes a
    spectral radius above one, and Koopman~(Stable) forces every eigenvalue onto
    the unit circle through a Cayley transform. The inset sketches the error
    growth each implies.}
    \label{fig:architecture}
\end{figure*}

\subsection{Four Dynamics Modules, One Encoder and One Objective}
\label{sec:variants}
% [orig 4 p1, p2] DONE, verbatim except that the two forward references now
% point at appendices B and C, where the derivations live.
All four models share the same encoder and JEPA training scheme~\citep{LeCun2022,Assran2023,Bardes2024}. The dynamics module is the only component that differs, and no pixels are decoded during training. An online encoder $\phi_\theta: \mathcal{X} \to \mathbb{R}^D$ maps each $128{\times}128$ RGB frame $\x_t$ to a latent vector $\z_t = \phi_\theta(\x_t)$ ($D{=}256$, selected by the tuning protocol of Appendix~\ref{app:impl}). A target encoder $\bar\phi$, parameterized identically to $\phi_\theta$, is updated as an exponential moving average (EMA) of the online weights~\citep{Grill2020} to produce stable prediction targets $\bar{\z}_{t+k} = \bar\phi(\x_{t+k})$. The dynamics module $f_\theta$ produces the $k$-step latent prediction $\hat{\z}_{t+k} = f_\theta^k(\z_t, \mathbf{a}_{t:t+k})$, conditioned on the executed action sequence. Figure~\ref{fig:architecture} shows the resulting architecture and the four variants.

For $f_\theta$ we compare four choices that separate nonlinearity, linearity, and spectral constraint. The \emph{MLP-JEPA baseline} uses a three-layer residual network, the standard nonlinear choice, subject to the exponential compounding of Equation~\ref{eq:compound}. The three \emph{Koopman} variants replace it with a single linear matrix $\K$. When actions are available the state evolves as $\z_{t+1} = \K\z_t + \B\mathbf{a}_t$, where $\B \in \mathbb{R}^{D \times m}$ is an unconstrained linear layer that maps $m$-dimensional actions into the latent space. The variants differ only in how the eigenvalues of $\K$ are constrained: \emph{Koopman~(Unconstrained)} optimizes $\K$ freely, \emph{Koopman~(Soft)} penalizes any spectral radius in excess of one (Appendix~\ref{app:soft}), and \emph{Koopman~(Stable)} enforces $|\lambda_i(\K)| = 1$ exactly through a Cayley parameterization (Appendix~\ref{app:dmd}). Initialization details, including the DMD warm-start used by the Stable variant, are given in Appendix~\ref{app:dmd}.

% [orig, training objective] DONE. Compressed from six sentences to four.
% "additionally" dropped as a banned filler word, and the decoder pointer moved
% from the deleted visual-MSE appendix to Section 5.2, which is where the one
% surviving decoder measurement is reported.
The prediction loss $\mathcal{L}_\text{pred}$ is the cosine distance between $\hat{\z}_{t+k}$ and $\bar{\z}_{t+k}$, averaged over all horizons up to $H$ (with $H{=}6$ for our training sequences of length 8). To prevent representational collapse, the online representations are regularized with VICReg~\citep{Bardes2022}. Its variance term $\mathcal{L}_\text{var}$ keeps each latent dimension informative and its covariance term $\mathcal{L}_\text{cov}$ decorrelates dimensions. The total base loss is
\begin{equation}
    \mathcal{L}_\text{base} = \lambda_p\,\mathcal{L}_\text{pred} + \lambda_v\bigl(\mathcal{L}_\text{var} + \lambda_c\,\mathcal{L}_\text{cov}\bigr).
    \label{eq:base_loss}
\end{equation}
The prediction weight $\lambda_p$ is held at zero for a short warm-up period so the encoder can establish a spatial representation under VICReg alone, and the warm-up length and the weights $\lambda_p$, $\lambda_v$, $\lambda_c$ are selected per method by the protocol of Appendix~\ref{app:impl} (Table~\ref{tab:hyperparams}). Training proceeds in two phases: the encoder and dynamics module are first optimized jointly under this objective, then both are frozen and a convolutional decoder is trained separately for visualization, which plays no role in the dynamics training signal or in any reported metric.

\subsection{Environments, Protocol and Tuning}
\label{sec:protocol}
% [orig, results preamble] DONE, verbatim. The three paragraphs that set up the
% comparison move here from the old Results section, where they preceded the
% first table.
We evaluate the three Koopman variants against the nonlinear MLP-JEPA baseline on the pendulum- and cartpole-swingup tasks from the DeepMind Control Suite~\citep{Tassa2018}. To test visual robustness, we extend the evaluation to three levels of distraction complexity (\textit{easy}, \textit{medium}, \textit{hard}) using the Distracting Control Suite~\citep{Stone2021}, which combines dynamic natural-scene background videos with camera-pose and color perturbations of increasing magnitude (Figure~\ref{fig:environments}). All final results are computed on a held-out test set, disjoint from both training and model selection, averaged over 10 random seeds, and reported as means with 95\% bootstrap confidence intervals.

The MLP-JEPA baseline shares the encoder architecture, training objective, and data with the Koopman variants, and its remaining (optimization and loss) hyperparameters are tuned independently by the same held-out selection protocol (Appendix~\ref{app:impl}), so that the latent dynamics module is the only structural difference across the comparison and the baseline is not disadvantaged by shared hyperparameters. This MLP-on-JEPA design is also the standard choice in the JEPA literature~\citep{Bardes2024,zhou2025dinowmworldmodelspretrained}. We do not compare against world models with different representation objectives (e.g., Dreamer~\citep{Hafner2025}, TD-MPC2~\citep{Hansen2023}), since that would confound the representation objective with the dynamics module. The MLP dynamics module carries roughly three times more parameters than the linear Koopman dynamics ($\approx\!198$k vs.\ $\approx\!66$k), so the Koopman variants' advantage is not attributable to greater capacity.


\subsection{Measurements and Controls}
\label{sec:measurements}
% [draft 4.4 + orig] DONE. The draft's 699 words cut to about 440. The probe
% paragraph merges the approved "Evaluation Metrics" text with the approved
% "Probe Choice" limitation, which is why that limitation is not in Section 7.
% Every definition puts the description before the name, per the guide, and
% this is the only site where the five terms are emphasized.
% MOVED: the Riccati paragraph to appendix D, which is also where the solver is
% named, since a package name does not belong in the main text.
Section~\ref{sec:bound_limits} argued that a bounded error curve is not by itself evidence of a learned model. This subsection defines the readouts and the controls that separate the two. All of them reuse the trained checkpoints unchanged, and the controls differ from the learned models only in the dynamics.

\paragraph{Probe state MSE.}
To check whether the learned latent states capture the underlying physical dynamics, we fit a \emph{probe}~\citep{alain2016understanding}, a two-layer MLP with 64 hidden units trained with the world-model weights frozen, mapping latent vectors to ground-truth physical coordinates. The two hidden layers absorb natural coordinate changes such as angle to $(\sin\theta, \cos\theta)$, which a linear probe would penalize independently of the quality of the representation, so the probe tests whether physical structure is recoverable from the latent rather than whether it is linearly aligned with it. Probe state MSE is the squared error between the true physical trajectory and the probe's prediction over a latent rollout, and probe $R^2$ measures the single-frame encoding, with $R^2 < 0$ signalling collapse. The metric is reported in raw coordinate units, which carry no intrinsic scale, so every table also carries the constant predictor defined below as a reference level.

\paragraph{A scale-invariant second readout.}
The probe is fitted on encoded observations and then applied to latents the model has iterated on its own, so it is used away from the distribution it was fitted on, and it responds to the length of the vector it receives as well as to its direction. We add a readout with neither property. Encoding a bank of held-out frames, matching each predicted latent to the nearest bank entry by cosine similarity, and reporting the squared error of the true physical state of the frame retrieved gives a score that reads direction only, never returns a state the environment did not visit, and requires no fitting. We call this \emph{retrieval}.

\paragraph{Trivial predictors.}
Three controls replace the dynamics with something that cannot have learned them, while keeping each model's own encoder and probe. Copying the current latent forward unchanged, $\hat{\z}_{t+k} = \z_t$, is the operator $\K = I$ and is non-expansive. We call it \emph{persistence}. Returning the mean physical state of the training set at every horizon is the constant predictor of Section~\ref{sec:bound_limits}. We call it \emph{mean state}, and whichever of the two is more accurate in a given cell is the \emph{trivial floor}, the level any useful model has to clear. A third control substitutes a frozen \emph{random orthogonal} matrix for $\K$, which satisfies the Cayley variant's spectral guarantee exactly while having learned nothing, and so separates the value of the constraint from the value of training under it. The first two depend on the encoder only through the probe, and vary by under one percent across the four, so we pool them per environment into one row each.

\paragraph{Spectral rescaling.}
Rescaling a learned operator to a target spectral radius, $\K \leftarrow (\rho_\text{target} / \rho(\K))\,\K$, and re-running the rollout with the encoder, the probe and the action matrix untouched varies the one quantity the theory identifies while holding the learned features fixed. We call this \emph{spectral rescaling}, and use $\rho_\text{target} \in \{0.9, 0.95, 1.0, 1.05, 1.1, 1.2\}$.

% =============================================================================
